\documentclass[french]{article}
\usepackage[utf8]{inputenc}
\usepackage[T1]{fontenc}
\usepackage{lmodern}
\usepackage{geometry}
\usepackage{graphicx}
\usepackage{babel}
\usepackage{amsmath}
\usepackage{amsthm}
\usepackage{amssymb}
\usepackage{hyperref}
\usepackage{stmaryrd}
\usepackage{enumitem}
\usepackage{tcolorbox}
\usepackage{dsfont}
\usepackage{multirow}
\usepackage{etoc}
\usepackage{titlesec}
\usepackage{esint}
\usepackage{cancel}
\usepackage{mathdots}

\setlist[itemize]{label=$\bullet$}


\renewcommand{\P}{\mathbb{P}}
\newcommand{\N}{\mathbb{N}}
\newcommand{\Z}{\mathbb{Z}}
\newcommand{\Q}{\mathbb{Q}}
\newcommand{\R}{\mathbb{R}}
\newcommand{\C}{\mathbb{C}}
\newcommand{\K}{\mathbb{K}}
\renewcommand{\L}{\mathbb{L}}
\newcommand{\U}{\mathbb{U}}
\newcommand{\st}{^*}
\newcommand{\tq}{\middle|}
\newcommand{\inv}{^{-1}}
\newcommand{\E}{\mathbb{E}}
\newcommand{\F}{\mathbb{F}}
\newcommand{\V}{\mathbb{V}}
\renewcommand{\ker}{\text{Ker}}
\newcommand{\im}{\text{Im}}
\newcommand{\card}{\text{Card}}
\newcommand{\Tr}{\text{Tr}}

\title{TD Polynômes \\ Exercice 34}
\date{}
\begin{document}
\maketitle

\section{Énoncé}
Soit $n$ et $k$ deux entiers naturels. On considère des réels $x_1,\ldots,x_n$ deux à deux distincts et des réels $y_1,\ldots,y_n$ quelconques. Pour $P\in\R_{k-1}[X]$, on note $E_P=\{i\in\llbracket1,n\rrbracket\ : \ P(x_i)=y_i\}$.
\begin{enumerate}
	\item Donner une condition nécessaire et suffisante sur $k$ et $n$ pour que, quel que soit $(y_1,\ldots,y_n)$, il existe $P\in\R_{k-1}[X]$ tel que $\card(E_P)=n$. Y a-t-il alors unicité de $P$ ?
	\item On suppose $k< n$.
	\begin{enumerate}[label=\alph*.]
		\item Montrer qu'il existe au plus un polynôme $P\in\R_{k-1}[X]$ tel que $\displaystyle\card(E_P)\geqslant\frac{n+k}{2}$.
		\item Montrer qu'il existe deux polynômes $Q$ et $R$, avec $R\ne 0$, tels que $$\deg(Q)<\frac{n+k}{2},\ \deg(R)\leqslant\frac{n-k}{2}\ \text{et}\ \forall i\in\llbracket1,n\rrbracket,\ Q(x_i)-y_iR(x_i)=0$$
		\item Montrer que s'il existe $P\in\R_{k-1}[X]$ tel que $\displaystyle\card(E_P)\geqslant\frac{n+k}{2}$, alors $P=Q/R$.
	\end{enumerate}
\end{enumerate}
\section{Solution}

\end{document}
